\texttt{systemd}\index{systemd} is het eerste proces (daemon) dat opgestart wordt op een Linux systeem. Het is daarmee de moeder van alle processen. Als er een daemon opgestart moet worden dan wordt dat gedaan via \texttt{systemd}. \texttt{systemd} is dus de plek waar ook in de gaten gehouden kan worden wat er opgestart is en wat niet. \texttt{systemd} kan ook monitoren of een opgestart proces verdwenen (gecrashed) is, en daarmee kan het een bijzondere inkijk geven in de status van een systeem. Door deze informatie te koppelen aan een proces dat dat dienst kan doen als log-server kan er veel informatie over processen verzameld worden. De systemd software kent dan ook zijn eigen logging server genaamd \texttt{systemd-jourald}. Journald slaat zijn logs niet op als platte tekst bestanden, maar als binaire files in de \texttt{journal} directory van \texttt{/var/log}.

